\subsection{把八条助记符落实为内存中的三十二个字节}

汇编形式便于人阅读，但处理器取回的是字节。若不把两者对应起来，“操作者已经把代码写入
RAM”仍然缺少可核对对象。下面为本章用到的操作码规定确切数值：

\begin{longtable}{|L{0.18\linewidth}|L{0.18\linewidth}|L{0.50\linewidth}|}
\hline
\rowcolor{BookTableHead}
操作码 & 指令 & 四个字节的解释 \\ \hline
\code{0x10} & \code{MOVI} & 操作码，目标寄存器，16 位有符号立即数低字节、高字节 \\ \hline
\code{0x11} & \code{ADD} & 操作码，目标兼第一源寄存器，第二源寄存器，保留字节 \\ \hline
\code{0x12} & \code{ADDI} & 操作码，目标寄存器，16 位有符号立即数低字节、高字节 \\ \hline
\code{0x20} & \code{LOAD} & 操作码，目标寄存器，基址寄存器，8 位有符号偏移 \\ \hline
\code{0x21} & \code{STORE} & 操作码，源寄存器，基址寄存器，8 位有符号偏移 \\ \hline
\code{0x31} & \code{BNEZ} & 操作码，测试寄存器，16 位相对指令位移低字节、高字节 \\ \hline
\code{0x40} & \code{CALL} & 操作码，24 位相对字节位移 \\ \hline
\code{0x41} & \code{RET} & 操作码，三个必须为零的保留字节 \\ \hline
\end{longtable}

寄存器编号直接使用 \code{0x00--0x07} 表示 \code{r0--r7}。16 位字段也采用低有效字节
在前的顺序。分支位移以“从顺序后继开始相差多少条 4 字节指令”计数，因此循环分支从
\code{0x00100014} 的顺序后继 \code{0x00100018} 回到 \code{0x00100004}，相差
\(-5\) 条指令。16 位补码 \(-5\) 为 \code{0xfffb}，在内存中写成
\code{fb ff}。

\subsection{任务代码的地址、字节和含义逐行对应}

\begin{longtable}{|L{0.20\linewidth}|L{0.27\linewidth}|L{0.39\linewidth}|}
\hline
\rowcolor{BookTableHead}
起始地址 & 四个实际字节 & 译码结果 \\ \hline
\code{0x00100000} & \code{10 04 00 00} & \code{MOVI r4,0} \\ \hline
\code{0x00100004} & \code{20 05 00 00} & \code{LOAD r5,[r0+0]} \\ \hline
\code{0x00100008} & \code{11 04 05 00} & \code{ADD r4,r4,r5} \\ \hline
\code{0x0010000c} & \code{12 00 04 00} & \code{ADDI r0,4} \\ \hline
\code{0x00100010} & \code{12 01 ff ff} & \code{ADDI r1,-1} \\ \hline
\code{0x00100014} & \code{31 01 fb ff} & \code{BNEZ r1,-5} \\ \hline
\code{0x00100018} & \code{21 04 02 00} & \code{STORE r4,[r2+0]} \\ \hline
\code{0x0010001c} & \code{41 00 00 00} & \code{RET} \\ \hline
\end{longtable}

这张表让三个容易混淆的地址分开了。指令地址是 PC 的值；装入和存储指令形成的数据地址
来自寄存器与偏移；编码中的寄存器号只是选择处理器内部状态，不是内存地址。例如地址
\code{0x00100004} 保存的第二个字节 \code{05} 表示目标寄存器 \code{r5}，并不表示要
访问物理地址 5。

\subsection{完整译码一条装入指令}

当 PC 为 \code{0x00100004} 时，处理器取回
\code{20 05 00 00}。译码器按固定字段作出以下决定：

\begin{enumerate}
\item 第一个字节 \code{0x20} 选择 \code{LOAD} 语义；
\item 第二个字节 \code{0x05} 选择写回目标 \code{r5}；
\item 第三个字节 \code{0x00} 选择基址 \code{r0}；
\item 第四个字节 \code{0x00} 符号扩展为偏移 0；
\item 读取 \code{r0=0x00102000}，形成有效地址 \code{0x00102000}；
\item 发起 4 字节读事务；
\item 响应成功后把返回值写入 \code{r5}，PC 增加 4。
\end{enumerate}

其中第 5 步只产生地址，第 6 步才与外部存储交互，第 7 步才提交寄存器。若事务返回错误，
\code{r5} 必须保持旧值，普通顺序 PC 也不能提交。否则软件会看到一条“既失败又部分完成”
的装入指令。

\begin{center}
\begin{tikzpicture}[node distance=10mm and 14mm]
\node[stage,minimum width=2.7cm] (bits) {指令字节\\\code{20 05 00 00}};
\node[stage,minimum width=2.8cm,right=of bits] (decode) {字段译码\\op、rd、rb、d};
\node[stage,minimum width=2.8cm,right=of decode] (addr) {地址形成\\\code{r0+0}};
\node[stage,minimum width=2.8cm,right=of addr] (read) {RAM 读响应\\\code{0x00000007}};
\node[stage,minimum width=2.8cm,below=13mm of addr] (commit) {提交\\\code{r5=7, PC+=4}};
\draw[flow] (bits) -- (decode);
\draw[flow] (decode) -- (addr);
\draw[flow] (addr) -- (read);
\draw[flow] (read.south) |- (commit.east);
\end{tikzpicture}
\end{center}
\diagramnote{译码、地址形成、外部读取和状态提交是连续但不同的阶段。任何阶段失败，都不能
假装后继阶段已经发生。}

\subsection{负数为什么没有负号字节}

\(-2\) 的 32 位补码由 \(2^{32}-2\) 得到，即 \code{0xfffffffe}。低有效字节先存，所以
RAM 中依次是 \code{fe ff ff ff}。装入指令不附带“有符号”标签；它把 32 个位原样放入
\code{r5}。加法器对补码位串执行模 \(2^{32}\) 加法，得到与有符号加法一致的低 32 位。

第一轮后 \code{r4=0x00000007}。第二轮执行加法：
\[
\texttt{0x00000007}+\texttt{0xfffffffe}
=\texttt{0x00000005}\pmod{2^{32}}.
\]
最高位之外的进位被丢弃，寄存器得到 5。只有当上层要判断溢出或比较大小时，才需要区分
有符号解释与无符号解释。本次输入之和未超出 32 位有符号范围。

\subsection{对齐要求来自一次取回的组织方式}

教学处理器要求 4 字节指令地址和 4 字节数据地址均为 4 的倍数。因此
\code{0x00100004} 和 \code{0x0010200c} 合法，而 \code{0x00102002} 上的 4 字节读取
返回 \code{ALIGNMENT}。这个约束简化了取指和 RAM 字节通道的组合。

对齐不是“地址看起来整齐”的排版规则。若处理器允许非对齐读取，它可能要发出两笔事务并
拼接结果；若中间一笔失败，还要规定整条指令是否提交。不同体系结构可以选择不同接口，
但软件必须遵守当前机器的明确契约。

\subsection{保留字段为什么也必须检查}

\code{RET} 的后三个字节目前没有语义，规范要求它们为零。译码器若忽略任意值，未来给
这些字段增加新含义时，旧映像可能被解释成另一种指令。要求发送端置零、处理器拒绝非零
保留字段，可以使格式扩展保持明确。

\begin{problemframe}
\textbf{复算点。} 若循环分支的两个位移字节误写为 \code{fc ff}，符号扩展结果为 \(-4\)。
顺序后继 \code{0x00100018} 加上 \(-4\times4\) 得到
\code{0x00100008}，循环会跳过 \code{LOAD}，反复累加旧 \code{r5}。校验和能检测传输
改变；若映像本来就这样编码，只有执行语义检查或测试才能发现逻辑错误。
\end{problemframe}

\subsection{同一片 RAM 为什么可以同时放指令、数据和栈}

RAM 只根据地址返回字节。处理器在取指阶段把返回的四个字节送到指令译码器；装入阶段把
返回字节送到通用寄存器；\code{RET} 则把返回字节送到 PC。目标位置没有自带类型，访问
路径决定解释方式。

\begin{longtable}{|L{0.22\linewidth}|L{0.27\linewidth}|L{0.37\linewidth}|}
\hline
\rowcolor{BookTableHead}
访问时的 PC/指令 & 读取的 RAM 地址 & 返回字节的解释 \\ \hline
PC 取指 & \code{0x00100008} & \code{ADD r4,r5} 的编码 \\ \hline
\code{LOAD} & \code{0x00102004} & 32 位输入整数 \(-2\) \\ \hline
\code{RET} & \code{0x001efffc} & 新 PC \code{0x00000300} \\ \hline
\end{longtable}

这也意味着错误地址可能改变解释。若 PC 被破坏为 \code{0x00102000}，处理器会尝试把
\code{07 00 00 00} 当作指令，而不会得到“这是数组”的提示。若存储指令写到代码区域，
后续取指会看到修改后的字节。当前机器没有执行权限和只读页，所有保护都只是合作约定。
